\documentclass[11pt]{article}
\usepackage[a4paper,margin=0.9in]{geometry}
\usepackage{amsmath,amssymb}
\usepackage[T1]{fontenc}
\usepackage[utf8]{inputenc}
\usepackage{lmodern}
\usepackage{microtype}
\usepackage{algpseudocode}
\usepackage{enumitem}
\usepackage{titlesec}

\setlength{\parindent}{0pt}
\setlength{\parskip}{0.3em}
\setlist[itemize]{leftmargin=*,topsep=0.18em,itemsep=0.1em,parsep=0pt,partopsep=0pt}
\setlist[enumerate]{topsep=0.18em,itemsep=0.1em,parsep=0pt,partopsep=0pt}
\setlength{\abovedisplayskip}{7pt plus 1pt minus 2pt}
\setlength{\belowdisplayskip}{7pt plus 1pt minus 2pt}
\setlength{\abovedisplayshortskip}{5pt plus 1pt minus 2pt}
\setlength{\belowdisplayshortskip}{5pt plus 1pt minus 2pt}
\titlespacing*{\section}{0pt}{1.5ex plus .4ex minus .2ex}{0.7ex}
\titlespacing*{\subsection}{0pt}{1.1ex plus .3ex minus .2ex}{0.5ex}
\titlespacing*{\paragraph}{0pt}{0.8ex plus .2ex minus .1ex}{0.5em}

\newcommand{\R}{\mathbb{R}}

\begin{document}
{\centering
\Large\bfseries Report on \emph{The Saddle Point Problem of Polynomials}\par
\vspace{0.15em}
\normalsize Xiangyu LUO\par
\normalsize March 2026\par
}
\vspace{0.1em}

\section{Introduction}
Let $X \subseteq \R^n$, $Y \subseteq \R^m$, and let $F(x,y)$ be a polynomial. A pair $(x^\ast,y^\ast) \in X \times Y$ is a saddle point if
\[
F(x^\ast,y) \leq F(x^\ast,y^\ast) \leq F(x,y^\ast)
\quad
\text{for all } x \in X,\ y \in Y.
\]
This means that $x^\ast$ minimizes $F(\cdot,y^\ast)$ over $X$, while $y^\ast$ maximizes $F(x^\ast,\cdot)$ over $Y$. From the definition one immediately gets
\[
\max_{y \in Y}\min_{x \in X} F(x,y) \leq F(x^\ast,y^\ast) \leq \min_{x \in X}\max_{y \in Y} F(x,y),
\]
and both inequalities become equalities when a saddle point exists. Saddle points therefore connect pointwise optimality with the global min-max value.

This viewpoint is useful in several settings. In zero-sum games, an equilibrium is characterized by a min-max relation of exactly this kind. In robust optimization, one compares a decision with its worst admissible perturbation. Similar saddle-point formulations also appear in primal-dual models from imaging, inverse problems, and constrained optimization.

Classical results rely on convexity in $x$, concavity in $y$, and compact convex feasible sets. Polynomial problems need not satisfy these hypotheses; saddle points may fail to exist, stationary points need not be global, and direct elimination is expensive. The computational goal is to find a saddle point when one exists and otherwise certify nonexistence.

\section{Problem and Background}
The feasible sets are basic semialgebraic sets:
\[
X=\{x \in \R^n : g_i(x)=0 \ (i \in E_1^X),\ g_i(x)\ge 0 \ (i \in E_2^X)\},
\]
\[
Y=\{y \in \R^m : h_j(y)=0 \ (j \in E_1^Y),\ h_j(y)\ge 0 \ (j \in E_2^Y)\}.
\]
For a saddle point $(x^\ast,y^\ast)$, the points $x^\ast$ and $y^\ast$ solve the two lower-level problems
\[
\min_{x \in X} F(x,y^\ast),
\qquad
\max_{y \in Y} F(x^\ast,y).
\]
Hence KKT-type conditions are natural. Under a constraint qualification, there exist multipliers $\lambda_i,\mu_j$ such that
\[
\nabla_x F(x^\ast,y^\ast)=\sum_i \lambda_i \nabla g_i(x^\ast),
\qquad
\nabla_y F(x^\ast,y^\ast)=\sum_j \mu_j \nabla h_j(y^\ast),
\]
together with complementarity on inequality constraints. Here $a\perp b$ means $ab=0$. These relations are necessary, but not sufficient: they identify critical points, not necessarily saddle points.

If the constraint tuples are nonsingular, the multipliers can be written as polynomial maps
\[
\lambda(x,y)=(\lambda_i(x,y))_{i\in E_1^X\cup E_2^X},
\qquad
\mu(x,y)=(\mu_j(x,y))_{j\in E_1^Y\cup E_2^Y},
\]
obtained from polynomial left inverses of the corresponding constraint matrices. Consequently, every saddle point must satisfy
\[
\nabla_x F(x,y)=\sum_i \lambda_i(x,y)\nabla g_i(x),
\qquad
\nabla_y F(x,y)=\sum_j \mu_j(x,y)\nabla h_j(y),
\]
as well as the corresponding complementarity conditions. This converts the saddle-point question into a family of polynomial optimization problems. 

The following concepts appear repeatedly in the later sections.
\begin{itemize}
\item \textbf{Semialgebraic set:} a set described by finitely many polynomial equalities and inequalities.
\item \textbf{KKT relations:} stationarity and complementarity conditions for constrained local optima.
\item \textbf{Nonsingularity:} an assumption ensuring that the multipliers can be written polynomially in the variables.
\item \textbf{Genericity:} a property that holds outside an exceptional algebraic set of coefficients.
\item \textbf{Moment/SOS relaxation:} a semidefinite relaxation built from moment and localizing matrices.
\item \textbf{Flat truncation:} a rank condition used to detect exactness and extract optimizers.
\end{itemize}

Lasserre's hierarchy is the computational tool used to solve these polynomial subproblems. It replaces them with a sequence of semidefinite relaxations; once one relaxation is exact, the true minimizers or maximizers can be extracted.

\section{Methods}
The procedure alternates between candidate generation and global verification. Polynomial criticality equations first produce candidate points. Each candidate is then tested by the two one-sided optimization problems from the saddle-point definition. If one test fails, the optimizer of the failed subproblem yields an exclusion inequality, and the refined problem is solved again.

\begingroup
\begin{algorithmic}[1]
\State \textbf{Input:} polynomial objective $F(x,y)$, semialgebraic constraint tuples $g,h$, and multiplier maps $\lambda(x,y),\mu(x,y)$
\State \textbf{Initialize:} start with no exclusion inequalities and no verified saddle points
\Loop
    \State build and solve the current upper-level polynomial problem
    \If{the upper-level problem is infeasible}
        \State \Return no saddle point exists
    \EndIf
    \State collect candidate points $(x^\ast,y^\ast)$
    \ForAll{candidates $(x^\ast,y^\ast)$}
        \State compute $\vartheta_1(y^\ast)=\min_{x\in X}F(x,y^\ast)$
        \State compute $\vartheta_2(x^\ast)=\max_{y\in Y}F(x^\ast,y)$
        \If{$\vartheta_1(y^\ast)=F(x^\ast,y^\ast)=\vartheta_2(x^\ast)$}
            \State record $(x^\ast,y^\ast)$ as a saddle point
        \Else
            \State construct a valid exclusion inequality from the failed verification step
            \State add this inequality to the upper-level problem
        \EndIf
    \EndFor
    \If{a nonempty set of saddle points has been recorded}
        \State \Return the recorded saddle point set
    \EndIf
\EndLoop
\end{algorithmic}
\endgroup

\paragraph{Step 1: upper-level candidate problem.}
The upper-level problem minimizes $F(x,y)$ over $X\times Y$ while also imposing the stationarity and complementarity relations derived from $\lambda(x,y)$ and $\mu(x,y)$, together with any exclusion inequalities produced in earlier iterations. The criticality equations are
\[
\nabla_x F(x,y)-\sum_i \lambda_i(x,y)\nabla g_i(x)=0,
\qquad
\nabla_y F(x,y)-\sum_j \mu_j(x,y)\nabla h_j(y)=0,
\]
together with
\[
0\leq \lambda_i(x,y)\perp g_i(x)\geq 0 \quad (i\in E_2^X), \qquad
0\geq \mu_j(x,y)\perp h_j(y)\geq 0 \quad (j\in E_2^Y).
\]
Solving this problem produces candidate points. At this stage they satisfy the polynomial
criticality system, but not yet the global saddle inequalities.

\paragraph{Step 2: lower-level verification.}
To decide whether a candidate $(x^\ast,y^\ast)$ is genuine, one solves
\[
\vartheta_1(y^\ast):=\min_{x \in X} F(x,y^\ast),
\qquad
\vartheta_2(x^\ast):=\max_{y \in Y} F(x^\ast,y).
\]
These are the two one-sided optimization problems in the definition of a saddle point. Therefore,
\[
(x^\ast,y^\ast)\text{ is a saddle point }
\Longleftrightarrow
\vartheta_1(y^\ast)=F(x^\ast,y^\ast)=\vartheta_2(x^\ast).
\]
It is convenient to introduce the \emph{saddle gap}
\[
\gamma(x^\ast,y^\ast):=\vartheta_2(x^\ast)-\vartheta_1(y^\ast).
\]
By definition of $\vartheta_1$ and $\vartheta_2$ one always has
\[
\vartheta_1(y^\ast)\leq F(x^\ast,y^\ast)\leq \vartheta_2(x^\ast),
\]
so $\gamma(x^\ast,y^\ast)\geq 0$, and $\gamma(x^\ast,y^\ast)=0$ holds exactly when the candidate
satisfies the saddle inequalities. This step therefore supplies the global information missing from
the KKT system.

\paragraph{Step 3: exclusion inequalities.}
If verification fails, the optimizer of the failed lower-level problem yields a valid inequality
that removes the current bad candidate while preserving every true saddle point.

Suppose $u \in X$ satisfies $F(u,y^\ast)<F(x^\ast,y^\ast)$. Then $(x^\ast,y^\ast)$ is not a saddle point.
For any true saddle point $(x,y)$, the point $x$ minimizes $F(\cdot,y)$ over $X$, so
\[
F(x,y)\leq F(u,y),
\qquad\text{hence}\qquad
F(u,y)-F(x,y)\geq 0.
\]
Thus $F(u,y)-F(x,y)\ge 0$ is valid for all saddle points, but it is violated by the current false
candidate.

Likewise, if $v \in Y$ satisfies $F(x^\ast,v)>F(x^\ast,y^\ast)$, then every true saddle point $(x,y)$
satisfies
\[
F(x,v)\leq F(x,y),
\qquad\text{so}\qquad
F(x,v)-F(x,y)\leq 0.
\]
This gives the second type of update. Repeating the process either verifies a saddle point or
proves that the refined upper-level problem is infeasible.

\paragraph{Step 4: semidefinite realization.}
Each subproblem is a polynomial optimization problem of the form $\min p(z)$ subject to
$\phi(z)=0$ and $\psi(z)\ge 0$. Lasserre's hierarchy replaces it with a moment relaxation:
\[
\min \langle p,w\rangle
\quad
\text{subject to } M_k(w)\succeq 0,\ L_{\phi}^{(k)}(w)=0,\ L_{\psi}^{(k)}(w)\succeq 0.
\]
Here $w$ is a truncated moment sequence, $M_k(w)$ is the moment matrix, and
$L_{\phi}^{(k)}(w),L_{\psi}^{(k)}(w)$ are localizing matrices encoding the equality and inequality
constraints. As the order $k$ increases, the relaxation becomes tighter. When flat truncation
holds, the optimizer can be extracted from the moment data; if a sufficiently high relaxation is
infeasible, then so is the original polynomial problem.

The archimedean-type assumptions used in the convergence statements are important here. They
prevent feasible moment sequences from escaping to infinity and make it possible for the relaxation hierarchy to recover
actual minimizers and maximizers after finitely many steps.

\section{Results and Discussion}
The paper treats saddle-point problems for a polynomial objective over basic semialgebraic sets.
Its central idea is to represent the Lagrange multipliers by polynomial maps $\lambda(x,y)$ and
$\mu(x,y)$ when the constraint tuples are nonsingular. This converts the search for saddle points
into an iterative polynomial optimization procedure: one first solves an upper-level system to obtain
critical candidates, then checks each candidate through the two global one-sided problems
$\min_{x\in X}F(x,y^\ast)$ and $\max_{y\in Y}F(x^\ast,y)$, and finally adds exclusion inequalities
whenever a candidate fails the saddle inequalities. The resulting subproblems are treated by
Lasserre-type moment relaxations, so the algorithm combines algebraic optimality conditions with
global semidefinite certification.

\subsection{Results}
Let $K_0$ denote the feasible set of the initial upper-level problem and let $S_a$ denote the
saddle-point set. The first termination statement is explicit: if $K_0\setminus S_a$ is finite, then
the refinement loop stops after at most $|K_0\setminus S_a|$ iterations. In the generic case, the
criticality system has only finitely many complex solutions, and Theorem 3.4 bounds their number by
a degree-dependent quantity $M$. This yields a uniform upper bound of $M$ on the number of
refinement steps and shows that the iteration count is controlled by the algebraic complexity of the
input data.

The semidefinite part of the analysis is also quantitative. For the upper-level problem and for the
two lower-level verification problems, the paper proves finite convergence of the corresponding
moment relaxations under its stated conditions. Once the relaxation order is sufficiently large, the
SDP either becomes infeasible, which certifies that the polynomial problem has no feasible point, or
satisfies flat truncation, which allows the true optimizer to be extracted from the moment data. The
numerical examples are consistent with this picture: several saddle-point instances terminate after
$1$ to $4$ iterations, one example requires $9$ iterations, and several nonexistence cases are
detected after $3$ or $4$ iterations. Reported running times range from a few seconds to about
$127$ seconds, while some direct KKT or quantifier-elimination computations in the comparison study
do not finish within six hours.

\subsection{Numerical illustration}
To test the procedure we implemented it in Python (\texttt{saddle\_point.py}, with the moment
relaxations in \texttt{polyopt.py}). Moment and localizing matrices are assembled as sparse linear
maps of the truncated moment vector and the SDPs are solved with Clarabel; dependent rows of the
localizing equalities are removed before solving, since interior-point solvers need full row rank.
Flatness is tested with a relative rank tolerance of $10^{-4}$ and minimizers are extracted with the
Henrion--Lasserre algorithm. The multiplier maps are those of the simplex,
$\lambda_0=x^\top\nabla_xF$, $\lambda_i=\partial_{x_i}F-x^\top\nabla_xF$, and of the interval
$\{1-x^2\ge0\}$, $\lambda=-xF'(x)/2$. Because the KKT set usually contains a continuum of points with
the same value of $F$, the upper-level objective is perturbed by a small generic linear term
($10^{-2}c^\top(x,y)$ with random $c$) so that the relaxation has a unique, flat minimizer; every
cut remains valid because it holds at all saddle points.

\begin{center}
\begin{tabular}{p{0.38\linewidth}p{0.54\linewidth}}
\hline
Instance & Output of the algorithm\\
\hline
Example 6.1(i): $X=Y=\Delta_3$, $F=x_1x_2+x_2x_3+x_3y_1+x_1y_3+y_1y_2+y_2y_3$ &
saddle point $x^\ast=(0,1,0)$, $y^\ast=(\tfrac12,\tfrac12,0)$, value $\tfrac14$, found in the first
round (saddle gap $4\cdot10^{-8}$)\\
$F=xy+x^2-y^2$ on $[-1,1]^2$ & saddle point $(0,0)$ in the first round\\
$F=(x-y)^2$ on $[-1,1]^2$ & two candidates fail ($\vartheta_2-\vartheta_1\approx2.7$ and $4.0$);
after two exclusion cuts the upper-level relaxation is infeasible, certifying that no saddle point
exists ($\min\max=1>0=\max\min$)\\
\hline
\end{tabular}
\end{center}

For Example 6.1(i) the computed $y^\ast$ differs from the point $(\tfrac14,\tfrac12,\tfrac14)$
reported in the paper, but both are saddle partners of $x^\ast$: $F(x^\ast,y)=y_2(1-y_2)$ is maximized by
every $y\in\Delta_3$ with $y_2=\tfrac12$, and for any such $y$,
$F(x,y)-\tfrac14=x_1x_2+x_2x_3+x_3y_1+x_1y_3\ge0$ on $\Delta_3$ with equality at $x^\ast$. The
saddle-point set is therefore the segment $\{x^\ast\}\times\{y\in\Delta_3:\ y_2=\tfrac12\}$, and the
algorithm returns one of its points. The whole run takes about 15 seconds.

\subsection{Discussion}
The main conceptual strength of the method is that it separates local algebraic structure from global
saddle-point verification. The polynomial system built from $\lambda(x,y)$ and $\mu(x,y)$ supplies a
finite candidate set under generic assumptions, but a candidate is accepted only after the two
global optimization problems confirm the saddle inequalities. In this sense, the procedure goes
beyond a stationary-point test: it does not merely list critical points, but distinguishes genuine
saddle points from false candidates. The exclusion inequalities are important for the same reason,
because each failed verification step contributes new information to the next upper-level problem
instead of restarting the computation from the beginning.

The limitations are tied to the same algebraic framework. The cleanest formulation depends on
nonsingularity, since otherwise the multiplier maps need not admit a uniform polynomial
representation in $(x,y)$, and the upper-level description becomes less direct. The convergence
results for the moment relaxations also rely on the assumptions stated in the paper, and the
computed optimizers remain numerical outputs of semidefinite solvers, so their accuracy is bounded
by solver tolerances. Moreover, the degree-based bound $M$ is mainly a worst-case estimate and does
not settle the overall complexity of the refinement process, especially in non-generic situations
where infinitely many non-saddle critical points may arise. These limitations also indicate the most
natural extensions: sharper iteration bounds, stronger treatment of degenerate cases, and broader
applications to polynomial zero-sum games, robust min-max models, and related semialgebraic
saddle-point problems.

\end{document}
